\begin{abstract}
Language models can adapt to new tasks by following instructions provided in context.
Reusing the internal states formed from these instructions could allow the resulting task capabilities to be shared across executions.
However, it remains unclear how the ability to transfer these states relates to the model’s continued dependence on the instructions.
We study this relationship using skills, sets of instructions that describe how to perform a task.
At different layers, we transfer states at the instruction positions from an execution using the correct skill into one using an unrelated skill.
Separately, we block attention to the skill in the original execution.
Across four models and two benchmarks, all eight comparisons show that transfer weakens at later layers while the original execution remains strongly dependent on the skill.
We then augment state transfer with key/value caches, the stored information used by attention.
At layer 16, this joint intervention improves accuracy over state transfer alone by 6.8 to 43.4 percentage points across all eight combinations.
For Qwen3-8B/MedCalc, accuracy rises from 3.5\% to 42.1\%, and blocking subsequent skill reading reduces it to 7.0\%.
ThComplete-source controls further show that donor skill-position KV alone does not recover donor performance; the depth dependence of joint repair varies across models and tasks.
ese findings establish a separation between using contextual skills and reusing their internal states, and demonstrate cache-assisted recovery across models and tasks.
\end{abstract}
